\section{Introduction}\label{sec:intro}

Most applications push background work through a queue: e-mail and push notifications, image resizing, search
re-indexing, webhook delivery, report exports. A single queue often carries jobs whose costs differ by orders of
magnitude. A report export that needs tens of seconds sits in front of a notification that needs a few milliseconds,
and under load the cheap job waits for the expensive one. This is head-of-line blocking, and it is the same mechanism
that makes tail latency hard in interactive services \citep{dean2013tail}. Mean and median latency suffer too, because
most jobs are the cheap ones.

\paragraph{What theory says to do.}
If each job's processing time were known, serving the shortest first would minimise mean response time on one server
\citep{schrage1968}. The result is old and robust. The same idea underlies web servers that favour small requests
\citep{harchol2003}, schedulers for data-processing clusters \citep{pastorelli2013}, and size-interval task
assignment across servers \citep{harchol1999sita}. It also has a known price: favouring short jobs delays long ones,
and the literature has spent considerable effort on whether and when that is unfair
\citep{bansal2001,wierman2003}. When sizes are unknown but their distribution is, the Gittins index is optimal
\citep{gittins1979,scully2020gittins}, and when nothing is known, the foreground--background policy favours jobs that
have received the least service \citep{nuyens2008}; GPU cluster managers use both ideas \citep{gu2019tiresias}.

\paragraph{What queues do today.}
Message queues largely do not use any of this. Amazon SQS offers standard and FIFO queues and no message priority
\citep{sqstypes}. Kafka is a log, and a partition is consumed in order \citep{kreps2011kafka}. RabbitMQ does support
per-message priorities \citep{rabbitmqpriority}, but the priority is a number the publisher must supply, and it only
reorders messages that are still queued: with consumer prefetch too high, messages are already handed out and
cannot be overtaken. So the burden falls on the operator, who splits queues or hand-assigns priorities per endpoint.
That requires knowing each endpoint's cost, goes stale when code or traffic changes, and can backfire: in our
benchmark a plausible misconfigured map was slower than doing nothing.

\paragraph{The gap.}
The scheduling literature mostly assumes that sizes are known, or that an estimate is given. The systems that
do estimate sizes work at other layers: per-job estimates inside a cluster scheduler \citep{pastorelli2013}, runtime
predictions for batch jobs in supercomputers \citep{tsafrir2007,gaussier2015}, or CPU scheduling inside a serverless
host \citep{fu2022sfs,fu2024alps}. A general-purpose message broker, whose handlers are arbitrary code, has no
declared job sizes, and no mechanism to learn them. Yet the broker sees something valuable: it delivers every message
and receives the acknowledgement, so it can measure how long each kind of message actually took.

\paragraph{Our approach.}
\sys{} closes that loop inside the broker. For each message key (the producer, optionally joined to a job-type
header) it maintains a decaying histogram of handler durations reported on acknowledgement. At ingress it looks up the
key's expected duration and maps it to one of a few priority levels, using boundaries taken from global quantiles. Cold,
stale or high-spread keys take a default level. No model is trained or served, and the operator configures no
priorities. Estimates will be wrong, and the literature says how: over-estimation is mild and under-estimation is
dangerous \citep{dellamico2014}. We therefore keep aging, so a long job misjudged as short, or one that has waited a
long time, cannot starve.

\paragraph{What we find.}
We evaluate with experiments whose configurations and pass criteria were frozen before the registered runs, and we
report the cases where the gate failed. In short: on a production-style workload the approach beats FIFO, hand-tuned
priorities and misconfigured priorities; the benefit replicates on separate cloud hardware and transfers to RabbitMQ
through a sidecar; and it breaks down in identifiable regimes: backlogs lasting hours (where aging erases the order),
workloads whose keys do not recur, and workloads where long jobs carry most of the work.

\paragraph{Contributions.}
\begin{itemize}
  \item A broker design in which per-key duration statistics, learned from acknowledgements, drive multi-level priority
        scheduling with no operator configuration; an opt-in worker-time weighting with pausing aging for sustained
        overload (Section~\ref{sec:design}).
  \item A pre-registered evaluation with an oracle reference, hand-tuned and misconfigured baselines, an AWS replication
        on separate hardware, and a comparison with RabbitMQ and Redis (Sections~\ref{sec:method}--\ref{sec:results}).
  \item Trace-driven simulation on four public traces from two serverless providers and two HPC logs, with the negative
        results kept and explained.
  \item A plain statement of the trade-off: for a non-preemptive queue, shortest-expected-first lowers median and mean
        latency but does not improve P99 when long jobs are common, and long jobs can lose.
\end{itemize}
All evidence is synthetic or trace-driven; we make no claim of a tail-latency gain or of production readiness
(Section~\ref{sec:limits}).
